\section{测试与代码审查}
%% ============================================================

\subsection{闭环验证流程}

不要仅仅让 Codex 做代码变更就结束。应当要求它完成完整的验证闭环：

\begin{enumerate}[nosep]
    \item 为变更编写或更新测试
    \item 运行相关的测试套件
    \item 检查 lint、格式化或类型检查
    \item 确认最终行为与需求一致
    \item 检查 diff 中是否存在 bug、回归或风险模式
\end{enumerate}

\begin{importantbox}{关键前提：让 Codex 知道什么是"好"}
Codex 能够执行这个验证闭环，但前提是它知道什么算"好"。这些信息可以来自 prompt 或 \texttt{AGENTS.md}，包括测试命令、lint 配置、代码风格期望等。
\end{importantbox}

\subsection{/review 命令}

\texttt{/review} 是一个内置的代码审查命令，支持多种审查模式：

\begin{itemize}
    \item 对比 base branch 进行 PR 风格审查
    \item 审查未提交的变更
    \item 审查某个特定 commit
    \item 使用自定义审查指令
\end{itemize}

如果团队有 \texttt{code\_review.md} 文件并在 \texttt{AGENTS.md} 中引用，Codex 在审查时也会遵循该指导。

\begin{knowledgebox}{GitHub Cloud 集成}
如果使用 GitHub Cloud，可以设置 Codex 自动审查 PR。在 OpenAI 内部，Codex 审查了 100\% 的 PR。可以启用自动审查，也可以在需要时通过 \texttt{@Codex} 触发审查。
\end{knowledgebox}

\subsection{本章小结}

Codex 不仅能生成代码，还能测试、检查和审查代码。通过 \texttt{/review} 命令和 \texttt{code\_review.md} 配置，实现团队级别一致的代码审查标准。在 GitHub Cloud 上还可以实现全自动 PR 审查。

%% ============================================================
